\chapter{绪论}

\section{研究背景与意义}

处理器微架构侧信道漏洞，近十年来一直是安全研究的活跃领域。从学术界到工业界，大量工作围绕该问题展开。

现代CPU为了指令吞吐率，在硬件层面引入了多种优化机制。多级缓存填补处理器与主存之间的速度鸿沟，分支预测器提前预测跳转方向，乱序执行引擎使指令在操作数就绪后立即发射，硬件预取器根据访问模式提前将数据预取至缓存。但这些微架构组件在提升性能的同时，也在执行路径上留下了可被外部观测的物理痕迹——访存延迟的细微涨落、功耗曲线的局部起伏、共享资源仲裁的顺序差异。探测方利用这些旁路信号，无需直接攻破权限隔离机制，即可推断出被攻击进程内部的敏感信息。

云计算的大范围部署，使微架构侧信道攻击已构成现实威胁。公有云平台上，成千上万的租户共享同一批物理服务器。虚拟机之间虽然逻辑隔离，底层CPU缓存、分支预测器、内存控制器等硬件资源仍然是共享的。Ristenpart等人早在2009年就证实了同一亚马逊EC2实例上并置虚拟机的可行性\cite{Ristenpart2009}，为后续跨虚拟机旁路攻击研究奠定了基础。Zhang等人进一步推进了该方向，在PaaS云环境中实现了跨租户的缓存侧信道攻击\cite{Zhang2014}。他们证明了：即使没有管理员权限，监听方仅依靠缓存访问模式即可推断同平台其他用户的敏感信息。国内学者同样关注到这一威胁——张伟娟等人对缓存侧信道攻防技术进行了系统综述\cite{ZhangWJ2023}，王崇等人从防御角度梳理了缓存侧信道的研究进展\cite{WangC2021}。

其根本原因在于：多租户共享硬件，该部署模型本身就是微架构侧信道攻击的温床。

AES是缓存旁路攻击的主要目标之一。就应用范围而言——TLS/SSL、磁盘加密、VPN隧道、数据库加密，几乎所有需要数据保密性的场景均采用AES\cite{NIST2001}。为了追求性能，AES的软件实现普遍将SubBytes、ShiftRows和MixColumns三个操作合并为预计算查找表，即T表\cite{Osvik2006}。查表时的内存访问模式直接受密钥值决定，缓存命中与未命中的时间差异即成为信息泄露的标志。再加上AES密钥扩展算法是确定且可逆的：探测方只要恢复最后一轮（和倒数第二轮）子密钥，即可逆推出原始密钥。

以AES的T表实现为例，每轮加密的核心运算均转化为查表：四张T表$Te_0$至$Te_3$各含256个32位入口，每张表1024字节，恰好占16个64字节的缓存行。查表索引由明文字节与轮密钥字节异或决定，索引的高4位决定访问哪个缓存行。探测方通过Flush+Reload监控T表缓存行的访问状态，样本累积足够多后，排除法可将每个密钥字节的256个候选逐步收敛至唯一值。

2018年Spectre\cite{Kocher2019}和Meltdown\cite{Lipp2018}的公开发布标志着微架构安全研究的新阶段。Spectre利用分支预测器的误预测诱导CPU进入瞬态执行窗口，在窗内访问受保护数据，再通过缓存侧信道将数据外泄。Meltdown则利用乱序执行绕过特权检查，直接读取内核地址空间。这意味着微架构优化本身已成为安全威胁的根源。Spectre和Meltdown的披露，使得缓存侧信道从一种独立的攻击路径，升级为更复杂的瞬态执行攻击链中的关键数据外泄环节——微架构侧信道（主要是缓存）是数据外泄的通道，而瞬态执行是通向该通道的门户。在此之后，Foreshadow\cite{Bulck2018}针对Intel SGX可信执行环境发起攻击。RIDL\cite{VanSchaik2019}和ZombieLoad\cite{Schwarz2019}分别利用处理器内部缓冲区和填充缓冲区泄露数据。Canella等人对瞬态执行攻击进行了系统分类\cite{Canella2019}，从该分类中可见微架构侧信道漏洞的层次。

另一方面，微架构侧信道漏洞的影响已超出学术圈。Intel、AMD、ARM等芯片厂商被迫进行微码更新、操作系统补丁、编译器缓解，每次均需在性能和安全性之间进行权衡，但并不能消除漏洞的影响。云服务商在多租户隔离方面必须考虑安全性，某些场景下甚至直接为安全敏感客户提供独占物理机。此外，软件供应链不断延伸，第三方库和开源组件广泛嵌入，平台需对侧信道暴露面进行持续度量和自动化巡检，对传统的渗透测试和安全审计流程进行微架构侧信道的漏洞排查。

本课题以微架构侧信道安全测试为主题，选取AES加密算法的Flush+Reload缓存侧信道攻击作为分析对象，设计并交付可运行的测试用例与评估报告，完成风险检测与效果度量，输出标准化的测试流程、实验数据，为配置加固、安全审计提供参考指标。

\section{国内外研究现状}

\subsection{缓存侧信道攻击研究}

侧信道攻击的思想可追溯到上世纪九十年代Kocher的工作。1996年，Kocher发现不同密钥值会使密码运算产生可测量的时间差，于是提出了计时攻击\cite{Kocher1996}。他证明了：运行时间本身即可泄露密钥信息。1999年，他又提出了差分功耗分析\cite{Kocher1999}，利用统计手段从功耗曲线中提取密钥相关信息。此后几乎所有侧信道攻击研究均可追溯至这两篇论文的思路：发现新的旁路信号、建立模型，进而逐步提升攻击效率。

2002年，Page从理论上分析了缓存能否作为密码分析的旁路\cite{Page2002}。结论是肯定的——缓存命中与未命中之间的时间差，可被用于推断密钥信息。2005年，Bernstein提出了利用缓存计时差异恢复AES密钥的完整方案\cite{Bernstein2005}，在真实系统上实现了远程缓存计时攻击。整个攻击无需物理接触目标机器，仅通过网络测量加密操作的响应时间，即可逐步缩小密钥搜索空间。同一时期，Percival的工作\cite{Percival2005}从另一个方向印证了该威胁：他针对OpenSSL的RSA实现，展示了缓存冲突如何导致密钥泄露。

2006年，Osvik、Shamir和Tromer给出了两种缓存攻击策略：Evict and Time和Prime and Probe。Evict and Time的思路是先填充满缓存将目标数据逐出，再触发受害方加密、测量时间。时间差异暴露了缓存访问模式。Prime and Probe则相反：先将目标缓存组填满自身数据（Prime阶段），待受害方执行完毕后观察哪些位置被替换（Probe阶段），从而推断受害方访问了哪些缓存组。Prime and Probe无需共享内存，适用范围更广，但精度受到缓存组粒度的限制。Tromer等人后来发表了更完整的分析\cite{Tromer2010}，将缓存攻击对AES的威胁进行了量化，即通过少量缓存访问观测即可恢复完整AES密钥。

Flush+Reload方法是Yarom和Falkner在2014年完整发表的\cite{Yarom2014}。高分辨率和跨核心能力使其迅速成为该领域的标志性工具。其核心方法是借助x86的CLFLUSH指令将指定缓存行从整个缓存层级中逐出，然后测量重载时间——如果目标数据被受害方带入缓存，再次访问的时间就短，否则就长。与Prime+Probe相比，Flush+Reload的分辨率达到缓存行级别（通常64字节），噪声也更低。其跨核心能力依赖于共享L3缓存和MESI一致性协议。Benger等人将其应用于椭圆曲线密码\cite{Benger2014}，证明仅凭少量侧信道信息即可恢复完整密钥。

Liu等人在2015年证明了LLC上Prime+Probe攻击的可行性\cite{Liu2015}。跨虚拟机场景下恢复AES密钥，该成果对云安全评估产生了深远影响。Gruss等人研究的Cache Template Attacks\cite{Gruss2015}将缓存攻击推向了自动化：无需预先获知受害方的内存布局，通过模板匹配即可自动发现和利用缓存的侧信道泄露。他们此后又提出了Flush+Flush攻击\cite{Gruss2017}，以CLFLUSH指令本身的执行时间代替内存访问时间作为信号，进一步降低了被检测的风险。

Shusterman等人在JavaScript等脚本环境中也进行了验证\cite{Shusterman2019}：浏览器沙箱内仍可实现缓存侧信道攻击。无需在目标机器上安装任何本地程序，网页中的一段JavaScript代码即可监视同机器上其他进程的缓存访问模式。这显著降低了攻击门槛。Gruss等人从另一个方向发力，利用JavaScript实现了Rowhammer攻击\cite{Gruss2016}，使浏览器环境对微架构安全的威胁引起了广泛关注。

将这十几年的演进串联来看，缓存侧信道攻击呈现出清晰的发展脉络，可划分为三个阶段。第一阶段为理论分析与概念验证（2002--2006年）：Page从理论上论证了缓存作为密码分析旁路的可行性，Bernstein在真实系统上实现了远程缓存计时攻击，Percival针对RSA实现验证了缓存冲突导致密钥泄露的风险，Osvik等人则确立了Evict+Time和Prime+Probe两种基本攻击范式。这一阶段的核心贡献在于证明了缓存侧信道威胁的现实性，并为后续研究奠定了方法论基础。第二阶段为实用化与精度提升（2007--2014年）：Flush+Reload方法的提出解决了高分辨率和跨核心两个关键问题，使跨核心、跨虚拟机的攻击从理论构想变为工程现实；Gullasch等人提出的消去法则大幅降低了对观测精度的要求，使噪声环境下的密钥恢复成为可能。该阶段的标志性特征是攻击从实验室环境走向真实云平台，Irazoqui等人在跨虚拟机场景下成功恢复AES密钥即为典型例证。第三阶段为自动化与攻击面扩展（2015年至今）：Cache Template Attacks使攻击发现实现了自动化，无需预先获知受害方内存布局即可自动识别和利用侧信道泄露；Shusterman等人在浏览器JavaScript环境中实现了缓存侧信道攻击，打破了攻击需在目标机器上运行本地代码的传统前提；更为重要的是，Spectre和Meltdown的披露使缓存侧信道从独立的攻击手段升级为瞬态执行攻击链中不可或缺的数据外泄环节，Gruss等人提出的Flush+Flush变体则进一步降低了攻击被检测的风险。总体而言，缓存侧信道攻击的发展趋势表现为：空间分辨率从缓存组级别提升至缓存行级别，攻击场景从本地扩展至跨核心、跨虚拟机乃至浏览器环境，攻击方法从手工构造演进为自动化模板匹配，侧信道角色从独立攻击路径演变为复杂攻击链的基础组件。

国内研究方面，学者们从多个角度对微架构侧信道安全进行了系统性的梳理与研究。吴晓慧等人对微架构瞬态执行攻击与防御方法进行了系统总结\cite{WuXH2020}，梳理了Spectre、Meltdown及其变种的技术原理和缓解方案。杨帆等人从可信执行环境角度梳理了软件侧信道攻击的研究进展\cite{YangF2023}，分析了SGX等可信执行环境面临侧信道攻击的脆弱性。唐博文等人从处理器时间侧信道角度综述了攻防技术的最新进展\cite{TangBW2024}，对时间侧信道的攻击模型和防御机制进行了分类比较。Ge等人的综述\cite{Ge2018}则从更宏观的视角对微架构计时攻击和防御措施进行了系统性整理，涵盖了从理论分析到工程实践的全链条。张伟娟等人对缓存侧信道攻防技术进行了系统综述，从攻击模型角度对现有攻击方法进行了分类比较。王崇等人从防御角度梳理了缓存侧信道的研究进展，对硬件级、操作系统级和软件级防御方案的效果和局限性进行了分析。国内研究的特点在于侧重系统性综述和防御技术梳理，在攻击方法的原创性实验和量化评估方面仍有深入空间。

\subsection{AES缓存攻击研究}

AES的T表实现是缓存旁路攻击的主要入口。标准AES的软件实现为提升性能，通常将SubBytes、ShiftRows和MixColumns三个操作合并为查表。每个输出字节由T表查表异或得到，查表索引直接受输入明文和轮密钥的影响。监听方通过观测T表访问模式，即可逐步缩小密钥候选空间。

Aciicmez和Koç在2006年提出了基于追踪驱动的AES缓存攻击\cite{Aciicmez2006b}。方法是将单次加密过程中所有缓存命中和未命中事件收集起来，构建一条完整的缓存访问轨迹，从中获取密钥信息。此后研究发现分支预测器亦可作为侧信道\cite{Aciicmez2007}，通过条件分支的执行路径推断AES内部状态。Neve和Seifert对访问驱动的缓存攻击进行了改进\cite{Neve2007}，提升了效率。

Gullasch等人在2011年将访问驱动缓存攻击推向实用\cite{Gullasch2011}。他们提出的消去法（Elimination Method）是AES缓存攻击领域的一个里程碑。核心思路是：探测方无需精确获知哪个缓存行被访问了，只需知道哪个缓存行未被访问。未命中的信息即足以将不可能的密钥候选排除。这一思路大幅降低了对观测精度的要求，在噪声较大的真实环境中仍可实现。本文的攻击框架亦将排除法作为核心方法之一。

跨虚拟机场景下的AES缓存攻击是云安全研究的重要课题。Irazoqui等人在2014年实现了跨虚拟机的AES缓存攻击\cite{Irazoqui2014}。他们利用LLC的共享特性，在一个虚拟机中监视另一个虚拟机中AES实现的缓存访问，数分钟即获取了完整密钥。2016年他们又将攻击应用到了跨处理器场景\cite{Irazoqui2016}——即使探测方和被攻击方运行于不同物理CPU上，只要共享一个LLC，攻击亦可成功。这些工作对云服务商的安全策略具有现实威胁。

从另一维度审视侧信道缓存攻击，按信息利用方式大致可划分为三类。

计时驱动，以Bernstein的成果为代表。测量加密操作的总执行时间，据此推断缓存访问模式。优势在于实现简单、无需共享内存，局限在于精度低、噪声多。

访问驱动，Osvik等人的Prime+Probe和Gullasch等人的消去法均属于此类。直接观测缓存行的访问状态获取信息，精度高，但前提是要么具备共享内存，要么对缓存布局有充分了解。

追踪驱动，Aciicmez的工作是典型。将单次加密的完整缓存访问轨迹完整记录后再进行分析，信息量最大，代价是观测难度亦最高。

本文采用的Flush+Reload属于访问驱动这一分支，配合排除法进行高精度密钥恢复。

综观AES缓存攻击的研究进展，可以观察到以下演进趋势：在攻击精度方面，从Bernstein的计时驱动攻击（仅能测量加密总时间）到Gullasch等人的消去法（利用缓存行级未命中信息），攻击的空间分辨率和信息利用率持续提升；在攻击场景方面，从单机本地攻击发展到Irazoqui等人的跨虚拟机攻击和跨处理器攻击，攻击的适用范围不断扩大；在密钥恢复策略方面，现有工作大多针对AES-128设计，对AES-192和AES-256的攻击方案相对匮乏，且鲜有研究探讨如何复用同一批监控数据实现多阶段密钥恢复。这些不足为本文的研究提供了切入点。

\subsection{微架构侧信道防御研究}

面对日益严峻的缓存侧信道威胁，学界和工业界从不同层面提出了应对方案。防御思路大致分布在硬件、操作系统和软件三个层级上。

硬件级防御的核心思想是修改缓存架构以消除侧信道泄露。Wang和Lee在2007年提出了新型缓存设计方案\cite{WangLee2007}，利用随机化缓存映射和分区来破坏监听方对缓存状态的观测能力。Kim等人提出的STEALTHMEM\cite{Kim2012}在操作系统层面进行缓存分区，为安全敏感进程划分专属缓存区域，从根本上阻断了共享缓存带来的侧信道风险。Intel此后推出的Cache Allocation Technology（CAT）亦采用类似路线，具备硬件分区能力，但粒度和灵活度均有限。Tromer等人对AES缓存攻击的防御进行了系统分析，涵盖查找表预加载、缓存行对齐等工程手段。

操作系统级防御主要依赖调度和隔离。核心绑定将敏感进程锁定在特定CPU核心上，避免与其他进程共享L1/L2缓存。虚拟机亲和性调度尽量将安全等级不同的虚拟机放置到不同物理CPU上。但其代价亦很明显：降低了调度灵活性，且无法对LLC级别攻击进行防御。

软件级防御部署最为便捷，亦是本文实验中重点评估的方向。常见的方法包括：敏感操作完成后主动刷新缓存行（CLFLUSH指令），确保缓存中不残留敏感数据痕迹；加密操作中插入随机延迟或噪声访存，干扰时间测量；常数时间实现，保证所有代码路径的执行时间和内存访问模式与密钥无关。常数时间实现从根本上解决问题，但需大量改造现有代码，且存在性能开销。

然而，各类防御方案在实际部署中各有局限性。硬件级方案需修改处理器设计，周期长、成本高，已部署的芯片无法修复。操作系统级方案依赖调度策略，资源利用率与安全性之间难以兼顾。软件级方案部署灵活，但效果参差不齐，这正是本研究的出发点：既关注攻击本身，亦关注防御措施的量化评估方法。

从防御研究的整体格局来看，现有工作存在以下不足：第一，大多数防御方案的效果评估停留在定性描述层面，缺乏统一的量化指标体系来客观度量防御措施对信息泄露的抑制程度；第二，针对AES-192和AES-256等较长密钥变体的攻击与防御评估较少，已有工作主要聚焦于AES-128；第三，噪声注入等概率性防御方案的有效性分析多基于理论推导，缺乏在真实攻击框架下的系统性实验验证。本文旨在填补上述空白，通过构建统一的攻击测试框架，建立包含泄露带宽在内的量化评估指标体系，对多种防御措施进行系统性的实验评估。

\subsection{瞬态执行攻击研究}

2018年初，Spectre和Meltdown公开发布，Spectre利用分支预测器的误预测诱导瞬态执行，Meltdown利用乱序执行绕过特权检查，二者均将缓存侧信道作为数据外泄通道。此后Foreshadow针对Intel SGX可信执行环境发起攻击，利用L1终端故障（L1 Terminal Fault）绕过SGX的内存隔离保护。RIDL和ZombieLoad分别利用处理器内部缓冲区和填充缓冲区泄露数据，表明即使不共享缓存行，微架构内部缓冲区亦可成为侧信道泄露源。Brasser等人\cite{Brasser2017}在SGX环境中亦验证了缓存侧信道攻击的可行性。Canella等人对瞬态执行攻击进行了统一分类，从误预测类型、瞬态执行指令窗口和数据外泄方式三个维度建立了系统化的分类框架。

这些工作指向同一事实：缓存侧信道是几乎所有瞬态执行攻击的数据外泄基础。在瞬态执行攻击时代，缓存侧信道攻防研究仍然是最底层的基础问题。瞬态执行攻击的涌现也从侧面凸显了缓存侧信道安全测试的重要性：只有充分理解缓存侧信道的信息泄露机理并建立有效的量化评估方法，才能为瞬态执行攻击的防御提供底层支撑。

\subsection{现有研究不足与本文创新性}

综合上述国内外研究现状，现有工作在以下方面存在不足，本文针对这些不足提出了相应的创新性贡献。

第一，针对AES-192和AES-256的缓存侧信道攻击方案较为匮乏。现有Flush and Reload攻击研究主要聚焦于AES-128，对AES-192和AES-256的攻击仅停留在理论可行性讨论层面，缺乏完整的密钥恢复方案和实验验证。本文提出了基于等效块解密的两阶段密钥恢复方法：第一阶段通过排除法恢复最后一轮子密钥，第二阶段复用同一批Flush+Reload监控数据，结合等效解密中间值的代数推导直接恢复等效解密轮密钥$Kd_1$，经MixColumns变换得到加密倒数第二轮子密钥。该方法的核心创新在于利用AES等效解密结构将缓存访问信息转化为倒数第二轮密钥约束条件，且两个阶段共享同一批监控数据，无需重复采样。

第二，现有防御措施的评估缺乏统一的量化指标体系。大多数研究对防御效果的描述停留在定性层面，如"攻击成功率下降"或"信息泄露减少"，缺乏客观的量化度量。本文建立了涵盖命中率、泄露带宽和攻击成功率的量化评估指标体系，其中泄露带宽定义为单次观测互信息的4倍与采样频率的乘积，从信息泄露速率角度刻画侧信道威胁程度，为不同防御方案的横向比较提供了统一基准。

第三，现有攻击框架大多针对单一密钥长度设计，缺乏对多种AES变体的统一支持。本文构建了一套统一处理AES-128/192/256三种密钥长度的攻击测试框架，采用模块化设计，通过差异化的密钥恢复流程适配不同密钥长度，同时保持T表监控策略和排除法核心逻辑的一致性，为微架构侧信道安全测试提供了可复用的工具支撑。

\section{本文研究内容}

本课题以微架构侧信道安全测试为总目标，围绕AES加密算法的Flush+Reload缓存旁路攻击展开。研究工作划分为三个层次：理论、设计实现、评估。

理论层，研究Flush+Reload攻击的底层机制和数学基础。具体包括：CPU缓存层次结构与一致性协议、共享内存与页面去重机制、高精度时间戳测量；AES分组密码的算法细节与T表加速实现；排除法的数学构造——从单样本分析到多样本累积的统计收敛，以及期望样本量的理论推导；针对AES-192和AES-256的两阶段密钥恢复策略：如何复用同一批监控数据，结合解密代数推导恢复倒数第二轮子密钥。理论层的工作目标是建立完整的攻击数学模型，为后续的设计实现提供推导依据。

设计实现层，构建了一套能统一处理AES-128/192/256三种密钥长度的攻击测试框架。框架采用父子进程模型：父进程作为监听方运行Flush+Reload采样，子进程作为被攻击方运行AES加密，进程间通过管道同步通信。共享内存模块利用POSIX共享内存机制存放T表数据，以缓存行为单位监视T表访问。密钥恢复模块实现了排除法，通过累积未命中信息逐步排除错误候选——AES-128恢复最后一轮子密钥后逆推原始密钥，AES-192和AES-256采用两阶段恢复流程：阶段一恢复最后一轮子密钥，阶段二复用同一批监控数据，结合等效解密中间值直接恢复等效解密轮密钥$Kd_1$，再经MixColumns变换得到加密倒数第二轮子密钥。量化评估模块定义了命中率、未命中率、泄露带宽、攻击成功率等指标，用于客观度量攻击效果和信息泄露规模。框架支持一键运行和批量测试。

评估层，在Linux平台上对随机噪声注入和缓存行刷新两种防御手段进行对照测试。首先获取默认系统配置下的基线数据，再逐一实验不同强度的缓解措施：CLFLUSH缓存刷新指令、不同概率的随机噪声注入等策略，记录统计指标变化并分析。通过量化对比不同防御措施下的命中率变化和泄露带宽衰减，评估各方案的实际防护效果和性能开销。三个层次各有分工：理论层推导数学公式，设计实现层编写代码并进行测试，评估层以实验数据得出量化结论。

\section{论文组织结构}

全文共五章。

第一章绪论。阐述研究背景与选题意义，围绕云多租户环境下的侧信道威胁，梳理Spectre、Meltdown等瞬态执行攻击对微架构安全研究格局的冲击。然后从缓存旁路攻击、AES缓存攻击、微架构侧信道防御、瞬态执行攻击四个方向综述国内外研究现状，点明本文的研究内容与目标。

第二章相关技术背景。阐述本工作依赖的关键技术：CPU缓存的分层结构与一致性协议、共享内存的实现机制、高精度时间戳的采集方式、AES分组密码的算法细节与T表加速实现，以及Flush+Reload攻击的基本原理和几种常见缓解措施。

第三章攻击方法设计与实现。详细阐述攻击方法的设计思路和工程实现：从排除法的数学构造出发，推导单样本分析和多样本累积的统计收敛过程。对AES-192和AES-256，详细推导两阶段密钥恢复——如何复用监控数据恢复倒数第二轮子密钥，再经密钥扩展逆运算反解原始密钥。然后介绍系统架构设计、共享内存实现、量化评估指标体系、缓解措施的工程实现。

第四章实验与分析。展示基准配置下AES-128/192/256三种密钥长度的攻击效果，给出命中率、密钥恢复成功率、泄露带宽等指标数据。在不同缓解措施（随机噪声注入、缓存行刷新）下进行对照测试，讨论性能开销与安全性之间的权衡，分析各防御方案的实际防护效果。

第五章总结与展望。总结全文研究工作，展望后续值得深入的方向。
